\documentclass[UTF8]{ctexart}
\usepackage{xeCJK}
\usepackage{geometry}
\geometry{left=2cm,right=2cm,top=2.5cm,bottom=2.5cm}
\setlength{\headheight}{12.7pt}

\usepackage{titlesec}
\titleformat{\section}
  {\normalfont\large\bfseries}
  {\thesection}
  {1em}
  {}
\titlespacing*{\section}
  {0pt}
  {3.5ex plus 1ex minus .2ex}
  {2.3ex plus .2ex}

\usepackage{amsmath}
\usepackage{amssymb}
\usepackage{amsthm}
\usepackage{enumitem}
\setlist[enumerate]{itemsep=0pt,topsep=0pt,parsep=0pt,leftmargin=0pt,itemindent=2.5em}
\setlist[itemize]{itemsep=0pt,topsep=0pt,parsep=0pt,leftmargin=0pt,itemindent=2.5em}
\usepackage{fancyhdr}
\pagestyle{fancy}
\fancyhf{}
\usepackage{graphicx}
\usepackage{hyperref}
\usepackage{indentfirst}
\usepackage{lmodern}


\begin{document}
\setlength{\abovedisplayskip}{1pt}
\setlength{\belowdisplayskip}{1pt}

\section{有限集}

\hypertarget{sec:定义1.1}{}
\noindent\textbf{定义1.1 (有限集)}

给定集合$a$, 如果$a\prec\omega$, 就称$a$是\textbf{有限集}, 否则称$a$是\textbf{无限集}.

\vspace{15pt}

\hypertarget{sec:定理1.1}{}
\noindent\textbf{定理1.1}

自然数都是有限集.

\noindent{\kaishu 证明.}

任取自然数$n$, 显然$n\preccurlyeq\omega$.
假设$n\approx\omega$, 取双射$f:n\to\omega$.

\vspace{3pt}
下面对$m\leqslant n$归纳证明$\exists M\in\omega\,\,\forall k<m:f(k)\leqslant M$.

\begin{enumerate}
    \item 首先显然$\exists M\in\omega\,\,\forall k<0:f(k)\leqslant M$;
    
    \item 其次, 任取自然数$m$, 假设存在自然数$M$使得$\forall k<m:f(k)\leqslant M$,
    再假设$m+1\leqslant n$.
    
    \hspace{2.17em}
    那么对任意的$k<m+1$有:
    \vspace{3pt}
    \[\left\{\begin{aligned}
        &k<m\quad\implies\quad f(k)\leqslant M\leqslant M+f(m),\\[-3pt]
        &k=m\quad\implies\quad f(k)=f(m)\leqslant M+f(m),
    \end{aligned}\right.\\[3pt]\]
    所以$\forall k<m+1:f(k)\leqslant M+f(m)$. 归纳完毕.
\end{enumerate}

\vspace{3pt}

从而存在自然数$M$使得$\forall k<n:f(k)\leqslant M$, 于是
$M+1=f(f^{-1}(M+1))\leqslant M$, 矛盾. \hfill $\qedsymbol$

\vspace{15pt}

鸽笼原理是有限集最重要的性质.

\vspace{10pt}

\hypertarget{sec:定理1.2}{}
\noindent\textbf{定理1.2 (鸽笼原理)}

有限集不可能与其
真子集\href{../../02 集合的一般构造/05 势关系#nameddest=sec:定义1.1}{等势}.

\noindent{\kaishu 证明.}

% 下面的证明我觉得很巧妙, 相比于直接归纳(从大双射构造小双射)的做法要自然得多.
% 只是这里有限集的定义略有不同.

假设有限集$a$与它的真子集$b$等势, 即存在双射$f:a\to b$.

取$x\in a\setminus b$, 然后递归定义$p:\omega\to a$ :
\[\forall n\in\omega:\quad p(n)=\left\{\begin{aligned}
    &x,\quad&&n=0,\\[-3pt]
    &f(p(n-1)),\quad&&n>0.
\end{aligned}\right.\\[3pt]\]
下证$p$是单射,
对$n$归纳证明$\forall k\in\omega:p(k)=p(n)\to k=n$.

首先对任意自然数$k$, 如果$p(k)=p(0)=x$, 假设$k\neq 0$则
\[
    x=p(k)=f(p(k-1))\in b,
\]
矛盾. 所以$k=0$.

其次任取自然数$n$, 假设$\forall k\in\omega:p(k)=p(n)\to k=n$.

此时对任意自然数$k$, 如果$p(k)=p(n+1)$, 那么$k\neq 0$, 从而
\[
    f(p(k-1))=f(p(n))\implies p(k-1)=p(n)\implies k-1=n\implies k=n+1.
\]

综上$p$是单射, 从而$\omega\preccurlyeq a$, 矛盾. \hfill $\qedsymbol$

% Winter spring summer fall
% You're the fairest of them all
% I love you

\end{document}